\section{Results}\label{sec:results}

Table~\ref{tab:rules} gives the verdicts: the diagnosis is shown on the Tennessee Eastman process
and not on the debutanizer, F1 fires, and no other rule decided on the data fires (the literature
rule F4 fired before the data; Section~\ref{sec:remedies:comp}). This section reports the cells
behind the verdicts; the supplementary material lists every cell and contrast.

\begin{table*}[t]
\centering\footnotesize
\caption{Seven cells at T1, $\theta = 4$: means over the 15 runs. Loop off: split conformal,
calibrated once. Median width: the run median of the finite interval widths, in mol\,\%; worst
rolling coverage: the minimum over 100 consecutive intervals; sd($\alpha$): the standard deviation
of the level; IACT of the misses. Last rows: the paired difference, compensated minus
uncompensated, under replay to $\varepsilon/2$, with its 95\,\% $t$ interval (coverages as
fractions).}
\label{tab:t1cells}
\input{\figdir/n1_table_t1_cells.tex}
\end{table*}

\begin{table*}[t]
\centering\footnotesize
\caption{The compensated loop (replay to $\varepsilon/2$) against four baseline configurations at
T1, $\theta = 4$: means over the 15 runs. Valid: mean $U$ at most 2\,\% and mean coverage within
2 points of 0.90. F7 width: the run median over all issued intervals, infinite ones included,
in mol\,\%.
Coverage-matched width change: each run's baseline width scaled to the loop's coverage in that run
by the ratio of Gaussian quantiles, against the loop's width; in brackets, the runs in which the
baseline is narrower. F7 fires if a valid baseline is at least 5\,\% narrower on average and
narrower in at least 12 of the 15 runs.}
\label{tab:baselines}
\input{\figdir/n1_table_baselines.tex}
\end{table*}

\subsection{The diagnosis on the Tennessee Eastman process}\label{sec:results:diagnosis}

With the step tuned on history that ignores the delay, the uninformative share rises with the
dead time (Fig.~\ref{fig:diagnosis}). In the primary cell, T1 at $\theta = 4$, 6.89\,\% of the
intervals were uninformative, 6.04 points more than at $\theta = 0$ on the same runs (97.5\,\%
CI 3.41 to 8.68; higher in 13 of 15 runs), while coverage read 0.8974 (8077 of 9000 labels;
Wilson 95\,\% interval 0.8910 to 0.9035, which contains 0.90): F2 is shown and F3 does not fire.
The rise grows with the dead time: 3.00 points at $\theta = 2$ and, in T2, 3.07 points at
$\theta = 4$ and 11.46 points at $\theta = 14$ (97.5\,\% CI 5.23 to 17.68), at coverages of 0.890 to
0.898; at T2 the pooled Wilson intervals lie just below 0.90 at every dead time, $\theta = 0$
included (pooled coverage 0.8903 to 0.8955). Compensation does not rescue a delay-blind step. At T2, $\theta = 14$,
the compensated loop issued 20.54\,\% uninformative intervals against 11.74\,\% without
compensation: the replay at $\theta = 0$ chose larger steps for the compensated residuals (0.010 to
0.015, against 0.001 to 0.010), and at $\theta = 14$ those steps issued more uninformative
intervals.

\begin{figure*}[t]
\centering
\includegraphics[width=\textwidth]{n1_speed_chart.pdf}
\caption{The analyzer-speed chart: width ratio, compensated over uncompensated, against the label
age $a = \theta + 1$ updates, in minutes for the Tennessee Eastman arms (left) and in samples for
the debutanizer (right). Open markers: predicted by Eq.~\eqref{eq:p3} from each run's calibration
residuals before the evaluation rows were read; filled markers: measured, as the ratio of the 0.90
quantiles of the evaluation scores, with 95\,\% $t$ intervals over the 15 runs (the debutanizer is
one stream). T3 updates once per analyzer value. Every measured ratio lies below its prediction.}
\label{fig:speed}
\end{figure*}

\subsection{Tuning}\label{sec:results:tuning}

Replay at the plant's dead time held the budget. At T1, $\theta = 4$, the uncompensated loop tuned
to $\varepsilon/2$ issued 0.30\,\% uninformative intervals (F5a; budget 2\,\%; two runs above
zero, at 3.67 and 0.83\,\%) at a coverage of 0.8926, and across the T1 and T2 cells the replay
kept the mean $U$ at or below 0.47\,\%, compensated or not (Fig.~\ref{fig:diagnosis}). Nine tunings, all at
T2's first operating point without compensation, were infeasible and ran at 0.001. The two
policies that each ignore one property of the loop did not hold the budget: against the replay,
delay-blind tuning added 6.59 points of $U$ (CI 4.18 to 9.00) and the white-noise rule 17.89
points (CI 15.04 to 20.74), both at nominal coverage. Their larger steps adapt faster locally,
with worst rolling coverage 5.0 and 8.3 points above the replay's 0.731 (Table~\ref{tab:t1cells}).
The white-noise model behind the rule failed as Section~\ref{sec:analysis:persistence} expects: it
underestimated $\sigma_\alpha$ by 33 to 73\,\% in all 47 uncompensated cells at the documented
dead times, and the infinite share by 81 to 100\,\% in the 24 cells that had infinite intervals
(F1 fires); with compensation it was off by more than 25\,\% in 40 of 41 cells. Fixed steps did not
hold the budget everywhere. The step of \citet{gibbs2021aci}, 0.005, held the budget at T1 but not at T2 (2.64\,\%
at $\theta = 4$ and 9.71\,\% at $\theta = 14$, uncompensated), and 0.05 gave 13 to 50\,\% in every
T1 and T2 cell with $\theta \ge 2$. Of the comparison loops tuned by the same replay, the interleaved chains, the
Smith-predictor form and decaying steps held the budget at T1, $\theta = 4$ ($U \le 0.48\,\%$,
coverage 0.891 to 0.897); projecting the level onto $[0.01, 0.5]$ removed the infinite intervals
but under-covered (0.858, and 0.841 with compensation).

\begin{figure*}[t]
\centering
\includegraphics[width=\textwidth]{n1_persistence.pdf}
\caption{Persistence and the step at T1, $\theta = 4$. (a) IACT of the open-loop misses ($\gamma =
0$, the window still rolling), per run, without and with compensation; grey lines join the runs,
large markers are means with 95\,\% $t$ intervals. (b) The step selected by replay to
$\varepsilon/2$ per operating point, without and with compensation.}
\label{fig:persistence}
\end{figure*}

\subsection{Compensation}\label{sec:results:comp}

Compensation narrowed the intervals at the same budget. At T1, $\theta = 4$, the compensated loop
tuned to $\varepsilon/2$ issued 0.18\,\% uninformative intervals (F5b; two runs above zero, at 0.83
and 1.83\,\%), and its run-median width was 0.461~mol\,\% smaller than without compensation (CI
0.200 to 0.723), 17.8\,\% of the uncompensated 2.586 (CI 7.7 to 28.0\,\%), at a coverage of 0.8961
against 0.8926 (Table~\ref{tab:t1cells}). Its worst rolling coverage was 6.53 points higher (CI
1.84 to 11.22; reported, not claimed, as Section~\ref{sec:remedies:comp} explains), and its misses
less persistent (IACT 3.97 against 7.84). On this plant, where
split conformal itself covered 0.900, the loop added nothing measurable to compensated split
conformal: $U$ $+0.18$ points, coverage $-0.37$ points (CI $-2.11$ to $+1.38$), width $+0.023$ (CI
$-0.102$ to $+0.149$) and worst rolling coverage $+2.53$ points (CI $-0.36$ to $+5.43$).
Section~\ref{sec:results:debutanizer} shows a stream where the loop raised coverage, by amounts
whose 95\,\% intervals just include 0. The negative control shows that the bookkeeping matters:
against stored pairing, arrival pairing (uncompensated, replay to $\varepsilon/2$) lowered coverage
by 0.64 points (CI 0.38 to 0.91) and the worst rolling coverage by 1.07 points (CI 0.03 to
2.10).

\begin{figure*}[t]
\centering
\includegraphics[width=\textwidth]{n1_debutanizer.pdf}
\caption{The debutanizer at $\theta = 3$: coverage over the last 100 evaluation rows for the
delayed loop (replay to $\varepsilon/2$) and for split conformal, (a) without and (b) with
compensation; the legend gives the coverage over all 719 evaluation rows. Grey line: the nominal
coverage, 0.90.}
\label{fig:debutanizer}
\end{figure*}

\subsection{Baselines}\label{sec:results:baselines}

No baseline was both valid and narrower (F7 does not fire; Table~\ref{tab:baselines}). All three
conformal PID configurations were valid. At coverage matched in each run, the matched scorecaster,
which shares the compensator, came within 1.5\,\% of the loop (narrower in 8 of 15 runs;
difference $-0.031$, CI $-0.093$ to $+0.031$); the quantile-regression scorecaster was 1.6\,\% wider
and the Theta scorecaster 17.7\,\% wider. SPCI under-covered (0.8673) and is not valid. The
delayed loop with the compensator thus matched the best conformal PID configuration with one
state and one tuned step, against three tuned constants. The PID-guided soft sensor of
\citet{meng2026}, whose update is a conformal PID variant without delay handling, is represented
here by conformal PID itself. On
the debutanizer at $\theta = 3$, the matched configuration covered 0.9013 at a width of 0.192,
against the loop's 0.8999 and 0.189.

\subsection{Predicting the width}\label{sec:results:p3}

Eq.~\eqref{eq:p3}, evaluated on each run's calibration residuals before any evaluation row was
read, predicted a mean width ratio of 0.880 at T1, $\theta = 4$; the ratio measured on the
evaluation rows was 0.835. The mean error over the runs, $-4.9$\,\% (CI $-13.2$ to $+3.5$\,\%),
lies inside the $\pm 10$\,\% of F6, although single runs erred from $-27.5$ to $+16.0$\,\%. In all
15 cells of the chart, the 12 drawn in Fig.~\ref{fig:speed} and the three of the gradient-boosting
arm, the measured ratio lies below the prediction on average, by 0.2 to 7.1\,\% in eight cells and
by 18 to 44\,\% in seven, at the shortest label ages (T1 and T2 at $\theta = 0$) and on the
debutanizer: on average, compensation bought at least the predicted narrowing. Read along the label age, the chart prices the analyzer: on the Tennessee Eastman
process, a label 3~min old instead of 15~min would let compensation cut the width to 0.47 of the
uncompensated width instead of 0.88 (predicted; measured 0.39 and 0.84).

\subsection{Persistence and the step}\label{sec:results:persistence}

F8's gate passed (mean predicted ratio 0.880, at most 0.95), and compensation made the misses less
persistent. The IACT of the open-loop misses ($\gamma = 0$, the window still rolling) fell by 4.69
(CI 1.55 to 7.82) and was lower in 12 of 15 runs (Fig.~\ref{fig:persistence}a); their lag-1
autocorrelation fell by 0.058 (CI 0.013 to 0.103) and the mean run of consecutive misses by 0.64
(CI 0.14 to 1.14), with the miss rate 1.0 point lower (CI $-0.2$ to $+2.2$). With compensation the
replay held the budget at a larger step, as claim~2$'$ states: 0.007 at all three operating points,
against 0.003, 0.003 and 0.007 without (larger at two, equal at the third;
Fig.~\ref{fig:persistence}b). The IACT fell in every arm, by 0.65 (T3) to 13.43 (T2,
$\theta = 0$), and on the debutanizer at $\theta = 3$ it fell from 24.7 to 5.2 while the replay's
step rose from 0.002 to 0.010.

\subsection{The debutanizer}\label{sec:results:debutanizer}

On the debutanizer the delay did not raise $U$ at the primary and secondary dead times: under
delay-blind tuning $U$ was 0 at $\theta = 0$, 1, 3 and 7, so F2 is not shown there and the
diagnosis is restricted to the simulated plant. At $\theta = 12$, the end of the documented range, $U$ rose by
5.42 points (97.5\,\% CI 0.00 to 18.36; 11 blocks), not shown either; F2's power on one stream of
this length was about 0.1 (Section~\ref{sec:study}). The stream showed something else. Split
conformal covered 0.7497, and its coverage over 100 rows fell to 0.06 late in the stream
(Fig.~\ref{fig:debutanizer}a), which suggests, without testing it, that the evaluation rows
differ from the training rows. Uncompensated, the delayed loop covered 0.8651 at $\theta = 3$,
11.5 points more than split conformal (CI $-0.1$ to $+27.3$), but stayed below 0.90 at every dead
time (0.844 to 0.878; 0.865 to 0.878 under delay-blind tuning). Compensated, it covered 0.8999 at a width of 0.189, against 0.8762 for
compensated split conformal ($+2.4$ points, CI $-0.7$ to $+5.3$), and 3.5 points more than without
compensation (CI $-3.8$ to $+11.4$) at 55\,\% of the width. The white-noise rule issued 22.53\,\%
uninformative intervals here.

\subsection{Robustness}\label{sec:results:robustness}

The diagnosis and the remedies held with gradient boosting as the point model: at T1, $\theta =
4$, F2's rise was 3.71 points (97.5\,\% CI 0.90 to 6.52), the replay held $U$ at 1.39\,\% without and
0.47\,\% with compensation, Eq.~\eqref{eq:p3} erred by $-0.2$\,\%, and compensation lowered the
IACT by 4.84. At T2's faster cadence the replay held $U$ at or below 0.47\,\%, and
Eq.~\eqref{eq:p3} erred by $-4.0$ and $-5.2$\,\% at $\theta = 4$ and 14 but by $-44.2$\,\% at
$\theta = 0$, on the safe side. At T3, which has no dead time, $U$ was 0.28 and 0.33\,\%,
Eq.~\eqref{eq:p3} erred by $-3.5$\,\%, and compensation still lowered the IACT, by 0.65 (CI 0.17 to
1.14), though its step was not larger at every operating point.
